% ch9_b 第9章 §9.1尾–§9.2 (书页 367–379 / p382–p394)
% 块边界判定：
%   起点——书页 367（p382）顶部为 Problem 9.1.3，新起条目，无 ch9_a 溢出内容，
%     故本块无 %JOIN%，也不设 %--A-TAIL-- 区（ch9_a 止于书页 366 末的问题 9.1.2）。
%   终点——书页 379（p394）止于 9.2.2 节首要点列表的前两条；第三条要点印在
%     书页 380（p395）顶部，连同其后的正文 "Just like the U(1) ..." 均属 ch9_c。
%     本块按"要点列表只收到本页末"处理：keypoints 环境在两条后闭合；
%     第三条的原文与建议译文见文件尾【块界备注】，请 ch9_c 以其文件的
%     %--B-TAIL-- 区或正文开头续排该条（环境需另行开启），勿弃勿重。

\subsection*{问题 9.1.3}

设 $|\Psi_{\text{mean}}\rangle$ 为 $\pi$ 通量拟设的平均场基态。计算 $\langle\Psi_{\text{mean}}|\pmb{S}_i\pmb{S}_{i+\pmb{y}}|\Psi_{\text{mean}}\rangle$。利用这一结果求变分平均场能量 $\langle\Psi_{\text{mean}}|H|\Psi_{\text{mean}}\rangle$，其中 $H$ 是具有近邻相互作用的海森堡模型（见式 (9.1.15)）。比较 $\pi$ 通量态与二聚体态的平均场能量。

\subsection*{问题 9.1.4}

\textbf{Dirac 费米子}

\begin{enumerate}
\item 证明：在连续极限下，$\pi$ 通量态中的低能自旋子由下式描述
\begin{equation*}
H=v\sum_{\pmb{k}\sim0}\lambda^{\dagger}_{\alpha,\pmb{k}}(k_x\Gamma_x+k_y\Gamma_y)\lambda_{\alpha,\pmb{k}}
\end{equation*}
其中 $\lambda^{\top}_{\alpha,\pmb{k}}=(f_{\alpha,\pmb{k}},f_{\alpha,\pmb{k}+\pmb{Q}})$，$\pmb{Q}=(\pi,\pi)$，且 $\Gamma_x^2=\Gamma_y^2=1$。求 $v$ 与 $\Gamma_{x,y}$。在实空间中，$H$ 可改写为 $H=\int\mathrm{d}^2x\,v\lambda^{\dagger}_{\alpha}(-\mathrm{i}\Gamma_x\partial_x-\mathrm{i}\Gamma_y\partial_y)\lambda_{\alpha}$。

\item 相应的拉格朗日量 $\mathcal{L}=\mathrm{i}\lambda^{\dagger}_{\alpha}\partial_t\lambda_{\alpha}-v\lambda^{\dagger}_{\alpha}(-\mathrm{i}\Gamma_x\partial_x-\mathrm{i}\Gamma_y\partial_y)\lambda_{\alpha}$ 可以改写为
\begin{equation*}
\mathcal{L}=\mathrm{i}\bar{\lambda}_{\alpha}(\gamma^0\partial_t+v\gamma^x\partial_x+v\gamma^y\partial_y)\lambda_{\alpha}
\end{equation*}
其中 $\bar{\lambda}_{\alpha}=\lambda^{\dagger}_{\alpha}\gamma^0$，$\gamma_0^2=1$，且 $(\Gamma_x,\Gamma_y)=(-\gamma^0\gamma^x,-\gamma^0\gamma^y)$。求 $\gamma^{0,x,y}$。由 $\mathcal{L}$ 描述的费米子称为无质量狄拉克费米子。

\item 证明 $(\gamma^0,\gamma^x,\gamma^y)$ 在 $1+2$ 维中满足下列狄拉克代数：
\begin{equation*}
\{\gamma^{\mu},\gamma^{\nu}\}=\eta^{\mu\nu},\qquad \mu,\nu=0,x,y,
\end{equation*}
其中 $\eta^{\mu\nu}$ 是对角矩阵，$\eta^{00}=1$ 且 $\eta^{xx}=\eta^{yy}=-1$。

\item 有质量的狄拉克费米子由如下拉格朗日量描述：
\begin{equation*}
\mathcal{L}=\mathrm{i}\bar{\lambda}_{\alpha}(\gamma^0\partial_t+v\gamma^x\partial_x+v\gamma^y\partial_y)\lambda_{\alpha}+m\bar{\lambda}_{\alpha}\lambda_{\alpha}.
\end{equation*}
求相应运动方程 $\mathrm{i}(\gamma^0\partial_t+v\gamma^x\partial_x+v\gamma^y\partial_y+m)\lambda_{\alpha}=0$ 的解。证明该费米子具有色散 $\omega_{\pmb{k}}=\sqrt{v^2\pmb{k}^2+m^2}$。

\item \emph{与 $U(1)$ 规范场的耦合}

用 3.7.1 节所述的程序，求最小耦合到一个 $U(1)$ 规范场的狄拉克费米子的拉格朗日量。确保所得的拉格朗日量具有 $U(1)$ 规范不变性。
\end{enumerate}

\subsection{如何消除无能隙 $U(1)$ 规范玻色子}

\begin{keypoints}
\item 如何得到稳定的平均场态：通过 Chern--Simons 项和/或 Anderson--Higgs 机制让规范涨落获得有限能隙。
\item 稳定的自旋液体总包含电中性的自旋 $1/2$ 自旋子，且自旋子之间只有短程相互作用。
\end{keypoints}

我们已经看到，在 $\pi$ 通量态中，无能隙 $U(1)$ 规范玻色子与无能隙费米子强烈相互作用，甚至直到零能。对于 $\pi$ 通量相，很难求出一阶平均场理论（即费米子--规范场耦合系统）的性质。一阶平均场理论虽然不会误导我们，但它似乎复杂得毫无用处。现在的问题是：能否构造一个稳定的平均场态，使平均场涨落在低能下相互作用微弱？如果可以，一阶平均场理论的性质就能容易地求得。

由于一阶平均场理论至少必须包含实施每格点单费米子约束的规范涨落，得到稳定平均场态的一个办法是让规范玻色子获得能隙。有能隙的规范玻色子只能传递费米子之间的短程相互作用。基于朗道费米液体理论的经验，我们知道如何处理具有短程相互作用的费米子。

首先让我们澄清"让规范玻色子获得能隙"是什么意思。我们知道，规范玻色子不过是平均场拟设 $\chi_{ij}$ 的涨落。这些涨落的动力学依赖于平均场拟设。因此，所谓"让规范玻色子获得能隙"，指的就是找到这样的平均场拟设，使拟设的集体模式涨落有能隙。

为给寻找稳定平均场拟设的探索提供动机，让我们先考虑 $U(1)$ 规范玻色子可以通过哪些方式获得能隙。事实上，在 3.7.5 节和 4.4 节中，我们已经遇到过 $U(1)$ 规范玻色子获得能隙的两种方式。第一种方式是 Anderson--Higgs 机制（Anderson--Higgs mechanism），即规范玻色子与凝聚的电荷玻色子耦合。耦合系统的动力学由
\begin{equation*}
\mathcal{L}=\frac{1}{8\pi g^2}(\mathbf{e}^2-b^2)+c_1a_0^2-c_2\pmb{a}^2
\end{equation*}
在 $2+1$ 维及连续极限下描述（见式 (3.7.17)），其中 $c_1$ 与 $c_2$ 来自凝聚的电荷玻色子。第二种方式是通过 Chern--Simons 项。如果平均场拟设使得满带具有非零的霍尔电导，那么 $U(1)$ 规范玻色子的动力学将由
\begin{equation*}
\mathcal{L}=\frac{1}{8\pi g^2}(\mathbf{e}^2-b^2)+\frac{K}{4\pi}\epsilon^{\mu\nu\lambda}a_{\mu}\partial_{\nu}a_{\lambda}.
\tag{9.1.21}
\end{equation*}
描述。在问题 9.1.5 中我们将求出，Chern--Simons 项给规范玻色子一个非零能隙。找到稳定平均场态的关键，在于找到实现上述两种机制之一的平均场拟设。

在这两种情形中，规范场都只能传递短程相互作用。其结果是，自旋子不被禁闭。自旋液体之上的准粒子由自由自旋子描述，它们携带自旋 $1/2$、电荷为零。在存在 Chern--Simons 项时，自旋子甚至可以具有分数统计。

%FIG{9.3}: {手征自旋态的平均场拟设。$\chi_{ij}$ 在箭头方向上为 $\mathrm{i}\chi$。}

\subsection*{问题 9.1.5}

证明由式 (9.1.21) 描述的涨落具有有限能隙。（可以通过计算经典运动方程来做到这一点。）

\subsection{手征自旋态}

\begin{keypoints}
\item 平均场手征自旋拟设给出一个平移和旋转对称的自旋液体态，但它破坏时间反演与宇称对称性。
\item 平均场手征自旋态是稳定的。由平均场理论得到的手征自旋液体确实代表真实自旋系统的一个稳定量子相。
\item 手征自旋液体含有分数化激发——自旋子。自旋子携带自旋 $1/2$、不带电荷，并具有分数统计。
\end{keypoints}

在这一节中，我们将讨论实现上述第二种机制的一个平均场态。该平均场态由拟设 $\chi_{ij}$ 描述，其中 $\chi_{ij}$ 是复数并产生通量。由式 (9.1.6) 描述的自旋子的行为就好像在磁场中运动。当通量与自旋子密度（即每格点一个费米子）有正确的公度关系（commensuration）时，整数个朗道能级（更确切地说，由于晶格的存在，是朗道能带）被完全填满。在这种情况下，有效格点规范理论因填满的朗道能级的有限霍尔电导而含有一个 Chern--Simons 项。我们将把这样的稳定平均场态称为手征自旋态（chiral spin state）（Wen et al., 1989; Khveshchenko and Wiegmann, 1989）。

最简单的手征自旋态由下列拟设给出（见图 9.3）：
\begin{equation*}
\begin{array}{lll}
\chi_{i,i+\pmb{x}}=\mathrm{i}\chi_1, & \chi_{i,i+\pmb{y}}=\mathrm{i}\chi_1(-)^{i_x}, & a_0(i)=0,\\[2pt]
\chi_{i,i+\pmb{x}+\pmb{y}}=-\mathrm{i}\chi_2(-)^{i_x}, & \chi_{i,i+\pmb{x}-\pmb{y}}=\mathrm{i}\chi_2(-)^{i_x} &
\end{array}
\tag{9.1.22}
\end{equation*}

上述 $\chi_{ij}$ 在每个正方形中产生 $\pi$ 通量，在每个三角形中产生 $\frac{1}{2}\pi$ 通量。上一节讨论的 $\pi$ 通量相有些特殊：由于 $\pi$ 通量等价于 $-\pi$ 通量，$\pi$ 通量相保持时间反演对称性（$T$）与宇称（$P$）。然而，由式 (9.1.22) 描述的 $\frac{1}{2}\pi$ 通量相并不等价于 $-\frac{1}{2}\pi$ 通量相。在 $T$ 或 $P$ 之下，$\frac{1}{2}\pi$ 通量变为 $-\frac{1}{2}\pi$ 通量，式 (9.1.22) 中的 $\chi_2$ 变为 $-\chi_2$。因此，具有非零 $\chi_2$ 的手征自旋态自发破坏 $T$ 与 $P$。利用恒等式
\begin{equation*}
E_{123}\equiv\pmb{S}_1\cdot(\pmb{S}_2\times\pmb{S}_3)=2\mathrm{i}(\hat{\chi}_{12}\hat{\chi}_{23}\hat{\chi}_{31}-\hat{\chi}_{13}\hat{\chi}_{32}\hat{\chi}_{21}),\qquad \hat{\chi}_{ij}=f^{\dagger}_if_j
\tag{9.1.23}
\end{equation*}
我们发现 $\langle E_{123}\rangle\sim\mathrm{Im}(\chi_{12}\chi_{23}\chi_{31})$ 在手征自旋态中不为零。注意 $E_{123}$ 在 $T$ 或 $P$ 下为奇；因此可以把 $E_{123}$ 视为 $T$ 与 $P$ 破坏的序参量。

指定了手征自旋态之后，我们想知道哪个自旋哈密顿量支持手征自旋态。让我们考虑下列阻挫自旋哈密顿量：
\begin{equation*}
H=J_1\sum_i(\pmb{S}_i\cdot\pmb{S}_{i+\pmb{x}}+\pmb{S}_i\cdot\pmb{S}_{i+\pmb{y}})+J_2\sum_i(\pmb{S}_i\cdot\pmb{S}_{i+\pmb{x}+\pmb{y}}+\pmb{S}_i\cdot\pmb{S}_{i+\pmb{x}-\pmb{y}})
\tag{9.1.24}
\end{equation*}

我们想知道平均场哈密顿量 (9.1.6) 何时支持平均场手征自旋态。在平均场哈密顿量 (9.1.6) 中，$J_{ij}$ 对近邻取 $J_1$，对次近邻取 $J_2$。

当 $\chi_2=0$ 时，由平均场哈密顿量决定的自旋子能谱由式 (9.1.19) 给出。导带与价带在 $(k_x,k_y)=(0,0)$ 与 $(0,\frac{\pi}{a})$ 两点相接触。当 $\chi_2\neq0$ 时，导带与价带之间打开一个能隙。平均场自旋子能谱由
\begin{equation*}
E^{\pm}_{\pmb{k}}=\pm\sqrt{J_1^2\chi_1^2(\sin(k_x)^2+\sin(k_y)^2)+J_2^2\chi_2^2(\cos(k_x+k_y)+\cos(k_x-k_y))^2}
\tag{9.1.25}
\end{equation*}
给出。平均场基态通过填满价带得到。由于存在能隙，自旋--自旋关联是短程的。利用基态波函数，我们可以计算 $\langle f^{\dagger}_{\alpha i}f_{\alpha i}\rangle$ 来检验自洽条件 (9.1.7)。我们发现，当 $J_2/J_1<0.49$ 时，自洽方程只支持一个 $\chi_1\neq0$、$\chi_2=0$ 的解，即 $\pi$ 通量相。当 $J_2/J_1>0.49$ 时，式 (9.1.7) 还支持一个 $\chi_2\neq0$ 的第二解。发现第二解具有更低的平均场能量。在这种情况下，$T$ 与 $P$ 被自发破坏。

注意，手征态的平均哈密顿量等价于磁场中电子跳跃的问题。从属费米子与规范场 $a_{\mu}$ 的耦合，与电子和电磁场的耦合完全相同。因此可以预期，由式 (9.1.10) 描述的从属费米子系统具有一种类似于霍尔效应的现象。这里的"霍尔效应"意味着 $a_{\mu}$ 规范场的"电"场在横向感应出一个自旋子流：
\begin{equation*}
j_x=\sigma_{xy}e_y,\qquad e_y=\partial_ta_y-\partial_ya_0
\tag{9.1.26}
\end{equation*}
其中 $\sigma_{xy}$ 是霍尔电导。有一个定理指出，满带的霍尔电导总是量子化为整数乘 $1/2\pi$（Thouless et al., 1982; Avron et al., 1983）。在我们的情形，这意味着 $\sigma_{xy}=2n/2\pi$，其中因子 2 来自自旋。"霍尔"电导的数值可以用 4.4 节讨论的方法求得。我们得到
\begin{equation*}
\sigma_{xy}=2/2\pi.
\tag{9.1.27}
\end{equation*}
理解这一结果的最简单方式，是注意到每一个自旋向上的自旋子对应一个通量量子，每一个自旋向下的自旋子也对应一个通量量子。这样，自旋向上与自旋向下的自旋子都具有填充因子 $\nu=1$。关掉晶格势之后，平均场手征自旋态中的价带就变成第一朗道能级。价带中的自旋向上与自旋向下的从属费米子各对"霍尔"电导贡献 $1/2\pi$。因此，总"霍尔"电导由式 (9.1.27) 给出。

规范涨落的有效作用量可以通过把式 (9.1.10) 中的自旋子积掉而得到。利用"霍尔"电导与 Chern--Simons 项之间的关系，我们可以容易地把连续极限下的有效作用量写成
\begin{equation*}
S=\int\mathrm{d}^3x\,\frac{1}{2}\sigma_{xy}a_{\mu}\partial_{\nu}a_{\lambda}\epsilon_{\mu\nu\lambda}+\frac{1}{8\pi g^2}(\mathbf{e}^2-v^2b^2)+\dots
\tag{9.1.28}
\end{equation*}
这里式 (9.1.28) 中的 $g^2$ 是自旋子能隙的量级，而 $v$ 是典型自旋子速度、为 $1/aJ$ 的量级。利用有效拉格朗日量，我们可以计算规范涨落的低能动力学性质。

由于非零的"霍尔"电导 $\sigma_{xy}$，我们可以\emph{不}在导带中产生自旋子、也不在价带中产生空穴而改变自旋子密度。这一性质对理解准粒子的性质十分重要。让我们缓慢地打开 $a_{\mu}$ 场的通量 $\Phi=\int\mathrm{d}^2x\,b$。如果通量分布在一个大区域，通量中的"磁"场 $b$ 就很小。此时价带与导带之间的能隙保持有限，价带中的所有能级都被一个自旋向上和一个自旋向下的自旋子填满。打开通量会感应出一个环形"电"场 $e_{\theta}$，由于非零的 $\sigma_{xy}$，它在径向 $\hat{\pmb{r}}$ 产生自旋子流。于是，一些电荷积累在原点附近。我们发现，被感应出的自旋子总数为
\begin{equation*}
N\equiv\sum_i(\langle f^{\dagger}_{\alpha i}f_{\alpha i}\rangle-1)=-\sigma_{xy}\Phi=-\frac{\Phi}{\pi}
\tag{9.1.29}
\end{equation*}

我们想强调，通量只改变自旋子密度，不感应任何自旋量子数。无论通量感应出多少自旋子，通量管携带的自旋总为零，因为价带中的每个能级都被一个自旋向上和一个自旋向下的自旋子填满。

现在我们可以讨论手征自旋态中准粒子激发的量子数了。平均场手征自旋态中最简单的激发可以通过往导带中加入一个自旋子得到。然而，这个激发不是物理的，因为加入导带的这个多余自旋子违反约束 (9.1.5)。为满足约束，我们可以加入通量来改变价带中的自旋子密度。由上面的讨论可以看到，导带中自旋子所产生的多余自旋子密度，可以通过引入 $\pi$ 通量来抵消（见式 (9.1.29)）。因此，手征自旋态中的物理准粒子是经 $\pi$ 通量\emph{缀饰}的自旋子。缀饰后的自旋子携带自旋 $1/2$，因为通量不能感应任何自旋量子数。然而，作为电荷与通量的束缚态（注意自旋子携带 $a_{\mu}$ 规范场的一个单位电荷），自旋子具有分数统计。交换两个缀饰自旋子等价于把一个自旋子绕另一个运动半周，这会产生相位 $\theta=\frac{1}{2}q\Phi$。这里 $q=1$ 是自旋子的电荷，$\Phi=\pi$ 是束缚在自旋子上的通量。于是，缀饰自旋子的统计角为 $\theta=\frac{\pi}{2}$。具有这种统计的粒子恰介于玻色子（$\theta=0$）与费米子（$\theta=\pi$）之间，称为半子（semion）。

为理解自旋子的低能动力学，我们想导出自旋子的有效拉格朗日量。先通过令 $a_{\mu}=0$ 忽略规范涨落。此时，导带或价带中的自旋子由式 (9.1.25) 给出的色散描述。这里 $E^+_{\pmb{k}}$（$E^-_{\pmb{k}}$）在 $(0,0)$ 与 $(0,\pi/a)$ 有两个极小（极大）。因此，在连续极限下我们有四种自旋子，由下列有效拉格朗日量描述：
\begin{equation*}
\mathcal{L}=\sum_{I=1,2}\left[-\mathrm{i}f^{\dagger}_{I\alpha}\partial_tf_{I\alpha}+\frac{1}{2m_s}f^{\dagger}_{I\alpha}\partial_i^2f_{I\alpha}-\mathrm{i}\bar{f}^{\dagger}_{I\alpha}\partial_t\bar{f}_{I\alpha}+\frac{1}{2m_s}\bar{f}^{\dagger}_{I\alpha}\partial_i^2\bar{f}_{I\alpha}\right]
\tag{9.1.30}
\end{equation*}
这里 $f_{1\alpha}$ 与 $f_{2\alpha}$ 对应导带两个极小附近的自旋子，$\bar{f}_{1\alpha}$ 与 $\bar{f}_{2\alpha}$ 对应价带极大附近的空穴。当 $a_{\mu}\neq0$ 时，自旋子与规范场的耦合可以把式 (9.1.30) 中的 $\partial_{\mu}$ 换成 $\partial_{\mu}\pm\mathrm{i}a_{\mu}$ 而得到。这种耦合形式由规范不变性的要求决定。把与规范场的耦合包括进来之后，总的有效拉格朗日量具有如下形式：
\begin{equation*}
\begin{aligned}
\mathcal{L}={}&\sum_{I=1,2}\left[-\mathrm{i}f^{\dagger}_{I\alpha}(\partial_t-\mathrm{i}a_0)f_{I\alpha}+\frac{1}{2m_s}f^{\dagger}_{I\alpha}(\partial_i-\mathrm{i}a_i)^2f_{I\alpha}\right]\\
&+\sum_{I=1,2}\left[-\mathrm{i}\bar{f}^{\dagger}_{I\alpha}(\partial_t+\mathrm{i}a_0)\bar{f}_{I\alpha}+\frac{1}{2m_s}\bar{f}^{\dagger}_{I\alpha}(\partial_i+\mathrm{i}a_i)^2\bar{f}_{I\alpha}\right]\\
&-\frac{2}{4\pi}a_{\mu}\partial_{\nu}a_{\lambda}\epsilon_{\mu\nu\lambda}+\frac{1}{8\pi g^2}(\mathbf{e}^2-v^2b^2)
\end{aligned}
\tag{9.1.31}
\end{equation*}
由运动方程 $\frac{\delta\mathcal{L}}{\delta a_0}=0$ 得到
\begin{equation*}
n_1+n_2=-\frac{1}{\pi}b
\tag{9.1.32}
\end{equation*}
其中 $n_I=\sum_{\alpha}(f^{\dagger}_{I\alpha}f_{I\alpha}-\bar{f}^{\dagger}_{I\alpha}\bar{f}_{I\alpha})$，$I=1,2$，是自旋子的密度。式 (9.1.32) 告诉我们，自旋子经 $\pi$ 通量缀饰，与先前的结果一致。自旋子的统计也可以直接由式 (9.1.31) 计算。

我们想指出，按照式 (9.1.31)，对每个固定的 $I$ 有两种准粒子，即导带中的自旋 $1/2$ 自旋子与价带中的自旋 $1/2$ 空穴。这个结果是\emph{不正确的}。对每个固定的 $I$ 值，只应有一种准粒子，即自旋 $1/2$ 自旋子。导带中的自旋子与价带中的空穴在投影之后给出同一个自旋子（Affleck et al., 1988; Dagotto et al., 1988）。在意识到手征自旋拟设实际上具有 $SU(2)$ 规范结构之后，这一重复计数问题即可解决（见问题 9.2.4）。

\subsection*{问题 9.1.6}

通过极小化平均场能量 (9.1.6)，求出手征自旋态拟设 (9.1.22) 中决定 $\chi_1$ 与 $\chi_2$ 的方程。

\subsection*{问题 9.1.7}

证明式 (9.1.25)。

\subsection*{问题 9.1.8}

利用 4.4 节的结果证明式 (9.1.27)。

\section{$SU(2)$ 投影构造}

\subsection{隐藏的 $SU(2)$ 规范结构}

\begin{keypoints}
\item 投影构造具有一个隐藏的 $SU(2)$ 规范结构。
\item 平均场拟设 $(U_{ij},a^l_0(\pmb{i}))$ 是物理自旋液体的多对一标记。$SU(2)$ 规范等价的拟设标记同一个物理波函数。
\item 物理波函数在一个对称变换下的不变性，只要求相应的平均场拟设在计及一个规范变换的意义下不变。
\end{keypoints}

\subsubsection{自旋子对凝聚}

在上一节中，我们用 Chern--Simons 项构造了一个稳定的平均场态。在这一节中，我们想考虑因 Anderson--Higgs 机制而稳定的另一类平均场态。为使 Anderson--Higgs 机制奏效，我们首先需要一个携带 $a_{\mu}$ 荷的玻色子。其次，带电玻色子应当具有适当的动力学，使它们能够凝聚。在由投影构造得到的平均场理论中，我们没有携带 $a_{\mu}$ 荷的玻色子，但有携带 $a_{\mu}$ 荷的费米子。于是，我们可以用带电费米子的对来构造带电玻色子，并让这些费米子对凝聚。我们看到，Anderson--Higgs 机制可以通过费米子对凝聚实现。\footnote{如果费米子是电子，那么费米子对凝聚态就是 BCS 超导态。}下面我们将把费米子对凝聚纳入我们的平均场拟设。我们希望这些平均场拟设代表稳定的平均场态。

记得用自旋子算符表示时，哈密顿量 (9.1.1) 可以改写为
\begin{equation*}
H=\sum_{\langle ij\rangle}-\frac{1}{2}J_{ij}\left(f^{\dagger}_{i\alpha}f_{j\alpha}f^{\dagger}_{j\beta}f_{i\beta}+\frac{1}{2}f^{\dagger}_{i\alpha}f_{i\alpha}f^{\dagger}_{j\beta}f_{j\beta}\right)
\tag{9.2.1}
\end{equation*}
其中我们加入了适当的常数项 $\sum_if^{\dagger}_{i\alpha}f_{i\alpha}$ 以得到上述结果。

为得到包含费米子对凝聚的平均场哈密顿量，我们把算符 $f^{\dagger}_{i\alpha}f_{j\beta}$ 与 $f_{i\alpha}f_{i\beta}$ 都换成它们的基态期望值
\begin{equation*}
\begin{array}{ll}
\eta_{ij}\epsilon_{\alpha\beta}=-2\langle f_{i\alpha}f_{j\beta}\rangle, & \eta_{ij}=\eta_{ji},\\
\chi_{ij}\delta_{\alpha\beta}=2\langle f^{\dagger}_{i\alpha}f_{j\beta}\rangle, & \chi_{ij}=\chi^{\dagger}_{ji}.
\end{array}
\tag{9.2.2}
\end{equation*}
$\chi_{ij}$ 项已包含在先前的平均场拟设中。$\eta_{ij}$ 项是新的，它描述费米子对凝聚。我们还把约束 (9.1.4) 换成其基态平均
\begin{equation*}
\langle f^{\dagger}_{i\alpha}f_{i\alpha}\rangle=1,\qquad \langle f_{i\alpha}f_{i\beta}\epsilon_{\alpha\beta}\rangle=0
\tag{9.2.3}
\end{equation*}
这一约束可以通过引入依赖格点且不依赖时间的拉格朗日乘子（Lagrangian multiplier）$a^l_0(\pmb{i})$，$l=1,2,3$，来实施。这样，我们得到下列包含费米子对凝聚的零阶平均场哈密顿量：
\begin{equation*}
\begin{aligned}
H_{\text{mean}}={}&\sum_{\langle ij\rangle}-\frac{3}{8}J_{ij}\left[(\chi_{ji}f^{\dagger}_{i\alpha}f_{j\alpha}+\eta_{ij}f^{\dagger}_{i\alpha}f^{\dagger}_{j\beta}\epsilon_{\alpha\beta}+h.c)-|\chi_{ij}|^2-|\eta_{ij}|^2\right]\\
&+\sum_i\left[a^3_0(f^{\dagger}_{i\alpha}f_{i\alpha}-1)+[(a^1_0+\mathrm{i}a^2_0)f_{i\alpha}f_{i\beta}\epsilon_{\alpha\beta}+h.c.]\right]
\end{aligned}
\tag{9.2.4}
\end{equation*}
这里，式 (9.1.6) 中的 $\chi_{ij}$ 与 $\eta_{ij}$ 必须满足自洽条件 (9.2.2)，而依赖格点的场 $a^l_0(\pmb{i})$ 的选取须使平均场基态满足式 (9.2.3)。这样的 $\chi_{ij}$、$\eta_{ij}$ 与 $a^l_0$ 给出一个平均场解。

对于具有近邻自旋耦合的海森堡模型（见式 (9.1.15)），下列带有费米子对凝聚的平均场拟设是平均场方程 (9.2.2) 的一个解：
\begin{equation*}
\begin{array}{ll}
\chi_{i,i+\pmb{x}}=\chi, & \chi_{i,i+\pmb{y}}=\chi,\\
\eta_{i,i+\pmb{x}}=\eta, & \eta_{i,i+\pmb{y}}=-\eta,\\
a^l_0=0. &
\end{array}
\tag{9.2.5}
\end{equation*}
这样的拟设实际上对应一个 \emph{d} 波 BCS 态。我们称之为 \emph{d} 波态。\footnote{我们注意到，费米子配对序参量 $\eta_{ij}$ 在 $90^{\circ}$ 旋转下变号。这类似于携带角动量 2 的波函数。这就是我们把这个态称为 \emph{d} 波态的原因。}由于费米子对凝聚，我们预期由上述拟设描述的平均场态因 Anderson--Higgs 机制而不含有无能隙 $U(1)$ 规范玻色子。

好，尽管一切显得非常契合、物理图像也显得非常合理，上述结果却被证明是不正确的。我们在哪里犯了错误？这里给出的物理图像与推理是正确的，错误原来出在数学上。我们曾宣称，平均场理论 (9.1.10)（或更一般地，式 (9.2.4)）具有 $U(1)$ 规范结构，使得由 $U(1)$ 规范变换 (9.1.13) 相联系的两个拟设在投影之后对应同一个物理自旋波函数（见式 (9.1.14)）。事实上，平均场理论具有更大的 $SU(2)$ 规范结构。由 $SU(2)$ 规范变换相联系的两个拟设对应同一个物理自旋波函数。不同的规范结构对我们理解平均场态的性质有深远的影响。特别是，$\chi_{ij}$、$\eta_{ij}$ 与 $a^l_0(\pmb{i})$ 的涨落描述 $SU(2)$ 规范涨落。

\subsubsection{$SU(2)$ 形式}

为理解平均场哈密顿量 (9.2.4) 与约束 (9.1.4) 中的 $SU(2)$ 规范结构（Affleck et al., 1988; Dagotto et al., 1988），我们引入二重态（doublet）
\begin{equation*}
\psi=\binom{\psi_1}{\psi_2}=\binom{f_{\uparrow}}{f^{\dagger}_{\downarrow}}
\tag{9.2.6}
\end{equation*}
和一个矩阵
\begin{equation*}
U_{ij}=\begin{pmatrix}\chi^{\dagger}_{ij} & \eta_{ij}\\ \eta^{\dagger}_{ij} & -\chi_{ij}\end{pmatrix}=U^{\dagger}_{ji}
\tag{9.2.7}
\end{equation*}
利用式 (9.2.6) 与 (9.2.7)，我们可以把式 (9.2.3) 与 (9.2.4) 改写为
\begin{equation*}
\langle\psi^{\dagger}_{\pmb{i}}\tau^l\psi_{\pmb{i}}\rangle=0,\qquad l=1,2,3
\tag{9.2.8}
\end{equation*}
\begin{equation*}
H_{\text{mean}}=\sum_{\langle ij\rangle}\frac{3}{8}J_{ij}\left[\frac{1}{2}\operatorname{Tr}(U^{\dagger}_{ij}U_{ij})-(\psi^{\dagger}_iU_{ij}\psi_j+h.c.)\right]+\sum_i a^l_0\psi^{\dagger}_i\tau^l\psi_i
\tag{9.2.9}
\end{equation*}
其中 $\tau^l$，$l=1,2,3$，是泡利矩阵。由式 (9.2.9) 可以清楚地看到，哈密顿量在局域 $SU(2)$ 变换 $W_i$ 下不变：
\begin{equation*}
\begin{array}{c}
\psi_i\rightarrow W_i\psi_i\\
U_{ij}\rightarrow W_iU_{ij}W^{\dagger}_j
\end{array}
\tag{9.2.10}
\end{equation*}

$SU(2)$ 规范结构实际上起源于式 (9.1.2)。这里的 $SU(2)$ 是使物理自旋算符保持不变的自旋子之间最一般的变换。由于物理哈密顿量是自旋算符的函数，这些变换就成为规范变换。

\subsubsection{$SU(2)$ 规范变换的含义}

与 $U(1)$ 规范变换一样，$SU(2)$ 规范变换有如下含义：由 $SU(2)$ 规范变换相联系的两个拟设对应同一个物理自旋波函数。为看到这一点，我们注意到，每个格点上的 $f$ 费米子态与 $\psi$ 费米子态有下列关系：
\begin{equation*}
\begin{array}{ll}
|0_f\rangle=\psi^{\dagger}_2|0_{\psi}\rangle, & f^{\dagger}_{\uparrow}f^{\dagger}_{\downarrow}|0_f\rangle=\psi^{\dagger}_1|0_{\psi}\rangle,\\
f^{\dagger}_{\downarrow}|0_f\rangle=|0_{\psi}\rangle, & f^{\dagger}_{\uparrow}|0_f\rangle=\psi^{\dagger}_1\psi^{\dagger}_2|0_{\psi}\rangle
\end{array}
\end{equation*}
其中 $|0_{\psi}\rangle$ 是没有 $\psi$ 费米子的态，即 $\psi_a|0_{\psi}\rangle=0$。于是，物理的每格点单 $f$ 费米子态对应每格点 $\psi$ 费米子数目为偶数的态。这些每格点偶 $\psi$ 费米子态是局域 $SU(2)$ 单态（即它们在每个格点上都是 $SU(2)$ 单态）。空的 $\psi$ 费米子态对应自旋向下，双占据的 $\psi$ 费米子态对应自旋向上。有了这一理解，我们就可以把平均场态投影到每格点偶 $\psi$ 费米子子空间，得到自旋系统的合法波函数。设 $\pmb{i}_1,\pmb{i}_2,\dots$ 为自旋向上的位置。物理自旋波函数可以写成 $\pmb{i}_1,\pmb{i}_2,\dots$ 的函数，即 $\Psi_{\text{spin}}(\pmb{i}_1,\pmb{i}_2,\dots)$。由于 $\pmb{i}_1,\pmb{i}_2,\dots$ 是仅有的具有两个 $\psi$ 费米子的格点，我们得到
\begin{equation*}
\Psi_{\text{spin}}(\{\pmb{i}_n\})=\langle 0_{\psi}|\prod_n\psi_{1,\pmb{i}_n}\psi_{2,\pmb{i}_n}|\Psi^{(U_{ij})}_{\text{mean}}\rangle.
\tag{9.2.11}
\end{equation*}
由于 $\langle 0_{\psi}|$ 与 $\psi_{1,\pmb{i}}\psi_{2,\pmb{i}}$ 在局域 $SU(2)$ 变换下不变，由 $SU(2)$ 规范变换 $U'_{ij}=W_iU_{ij}W^{\dagger}_j$ 相联系的两个拟设 $U_{ij}$ 与 $U'_{ij}$ 只是标记\emph{同一物理态}的两个不同标签：
\begin{equation*}
\langle 0_{\psi}|\prod_n\psi_{1,\pmb{i}_n}\psi_{2,\pmb{i}_n}|\Psi^{(U_{ij})}_{\text{mean}}\rangle=\langle 0_{\psi}|\prod_n\psi_{1,\pmb{i}_n}\psi_{2,\pmb{i}_n}|\Psi^{(W_iU_{ij}W^{\dagger}_j)}_{\text{mean}}\rangle
\tag{9.2.12}
\end{equation*}

平均场态与物理自旋波函数之间的关系 (9.2.11) 使我们能够从平均场拟设的变换构造出物理自旋波函数的变换。例如，具有平移后拟设 $U'_{ij}=U_{i-l,j-l}$ 的平均场态 $|\Psi^{(U'_{ij})}_{\text{mean}}\rangle$ 在投影之后产生平移了的物理自旋波函数。

显然，平移不变的拟设将导致平移不变的物理自旋波函数。然而，投影后物理波函数的平移对称性并不要求相应拟设的平移不变性。物理态是平移对称的，当且仅当平移后的拟设 $U'_{ij}$ 与原拟设 $U_{ij}$ 规范等价。我们看到，规范结构可能使我们对称性的分析复杂化，因为物理自旋波函数 $\Psi_{\text{spin}}(\{\alpha_i\})$ 可能比投影前的平均场态 $|\Psi^{(U_{ij})}_{\text{mean}}\rangle$ 具有更多的对称性。

\subsubsection{自旋旋转不变性}

我们注意到，$\psi$ 的两个分量都携带自旋向上。因此，自旋旋转对称性在我们的形式体系中不是显式的，很难判断式 (9.2.9) 是否描述一个自旋旋转不变的态。事实上，对满足 $U_{ij}=U^{\dagger}_{ji}$ 的一般 $U_{ij}$，式 (9.2.9) 未必描述自旋旋转不变的态。然而，如果 $U_{ij}$ 具有如下形式
\begin{equation*}
\begin{array}{l}
U_{ij}=\chi^{\mu}_{ij}\tau^{\mu},\qquad \mu=0,1,2,3,\\
\chi^0_{ij}=\text{imaginary},\qquad \chi^l_{ij}=\text{real},\qquad l=1,2,3,
\end{array}
\tag{9.2.13}
\end{equation*}
那么式 (9.2.9) 将描述一个自旋旋转不变的态。这是因为上述 $U_{ij}$ 可以改写成式 (9.2.7) 的形式；此时式 (9.2.9) 可以改写为式 (9.2.4)，其中自旋旋转不变性是显式的。在式 (9.2.13) 中，$\tau^0$ 是单位矩阵，$\tau^{1,2,3}$ 是泡利矩阵。

\subsubsection{变分方法}

我们想说明，除了自洽方程 (9.2.2) 之外，还有另一种求平均场解的办法。我们可以把 $H_{\text{mean}}$（见式 (9.2.9)）的平均场基态，即 $|\Psi^{(U_{ij})}_{\text{mean}}\rangle$，当作尝试波函数，把 $U_{ij}$ 当作变分参数。引入
\begin{equation*}
(\tilde{U}_{ij})_{\alpha\beta}\equiv-2\langle\Psi^{(U_{ij})}_{\text{mean}}|\psi_{i,\alpha}\psi^{\dagger}_{j,\beta}|\Psi^{(U_{ij})}_{\text{mean}}\rangle
\end{equation*}
我们发现 $|\Psi^{(U_{ij})}_{\text{mean}}\rangle$ 的平均场能量为
\begin{equation*}
E_{\text{mean}}(\{U_{ij}\})=-\sum_{\langle ij\rangle}\frac{3}{16}J_{ij}\operatorname{Tr}(\tilde{U}^{\dagger}_{ij}\tilde{U}_{ij})
\tag{9.2.14}
\end{equation*}
这里 $E_{\text{mean}}(\{U_{ij}\})$ 是 $U_{ij}$ 的泛函，具有 $SU(2)$ 规范不变性：
\begin{equation*}
E_{\text{mean}}(\{U_{ij}\})=E_{\text{mean}}(\{W_iU_{ij}W^{\dagger}_j\})
\end{equation*}
平均场解 $U_{ij}$ 现在可以通过极小化 $E_{\text{mean}}$ 得到。

\subsubsection{一阶平均场理论}

为得到一阶平均场理论，我们从由平均场哈密顿量
\begin{equation*}
H_{\text{mean}}=\sum_{\langle ij\rangle}-\frac{3}{8}J_{ij}(\psi^{\dagger}_i\bar{U}_{ij}\psi_j+h.c.)+\sum_i a^l_0\psi^{\dagger}_i\tau^l\psi_i
\tag{9.2.15}
\end{equation*}
描述的零阶平均场理论出发，其中 $\bar{U}_{ij}$ 是平均场解，满足自洽条件
\begin{equation*}
\bar{\chi}_{ij}=\langle\psi^{\dagger}_{i\alpha}\psi_{j\alpha}\rangle,\qquad \bar{\eta}_{ij}=-\langle\psi_{i\alpha}\psi_{i\beta}\epsilon_{\alpha\beta}\rangle
\tag{9.2.16}
\end{equation*}
（译注：式 (9.2.16) 第二个等式中两处下标原书均印作 $i$；按式 (9.2.2)，第二个下标应为 $j$。）

式 (9.2.15) 中的 $a^l_0$ 的选取须使式 (9.2.3) 得到满足。平均场基态附近的重要涨落是 $U_{ij}$ 的下列"相位"涨落：
\begin{equation*}
U_{ij}=\bar{U}_{ij}\mathrm{e}^{\mathrm{i}a_{ij}}
\tag{9.2.17}
\end{equation*}
其中 $a_{ij}=a^l_{ij}\tau^l$ 是一个 $2\times2$ 无迹厄米矩阵。由于 $SU(2)$ 规范结构，这些涨落就是 $SU(2)$ 规范涨落。为了在低能下得到自旋液体态定性正确的结果，我们必须把这些涨落包括进来。这引导我们走向一阶平均场理论
\begin{equation*}
H_{\text{mean}}=\sum_{\langle ij\rangle}-\frac{3}{8}J_{ij}(\psi^{\dagger}_i\bar{U}_{ij}\mathrm{e}^{\mathrm{i}a^l_{ij}\tau^l}\psi_j+h.c.)+\sum_i a^l_0\psi^{\dagger}_i\tau^l\psi_i
\tag{9.2.18}
\end{equation*}
它描述耦合到 $SU(2)$ 格点规范场的自旋子。

\begin{smallnote}
我们想指出，自旋液体的平均场拟设 $U_{ij}$ 可以分为两类：非阻挫拟设——$U_{ij}$ 只把偶格点连接到奇格点；和阻挫拟设——$U_{ij}$ 在两个偶格点之间和/或两个奇格点之间非零。非阻挫拟设的每个方格只有纯 $SU(2)$ 通量，而阻挫拟设在某些方格中，除 $SU(2)$ 通量之外还有 $\pi/2$ 整数倍的 $U(1)$ 通量。
\end{smallnote}

\subsection*{问题 9.2.1}

用式 (9.2.2) 由式 (9.2.1) 得到式 (9.2.4)。

\subsection*{问题 9.2.2}

$\pi$ 通量态与 \emph{d} 波态的等价性

\begin{enumerate}
\item 证明 $\pi$ 通量态 (9.1.18) 由 $SU(2)$ 链变量
\begin{equation*}
U_{i,i+\pmb{x}}=-\mathrm{i}\chi_1\tau^0,\qquad U_{i,i+\pmb{y}}=-\mathrm{i}\chi_1(-)^{i_x}\tau^0
\end{equation*}
描述。
\item 证明当 $\eta=\chi$ 时，\emph{d} 波态 (9.2.5) 由 $SU(2)$ 链变量
\begin{equation*}
U_{i,i+\pmb{x}}=\chi\tau^3+\chi\tau^1,\qquad U_{i,i+\pmb{y}}=\chi\tau^3-\chi\tau^1
\end{equation*}
描述。
\item 证明当 $\chi$ 与 $\chi_1$ 以某种方式相联系时，上述两个拟设规范等价。求把一个拟设变成另一个的 $SU(2)$ 规范变换 $W_i$。（提示：考虑 $SU(2)$ 规范变换 $(\mathrm{i}\tau^1)^{i_x}(\mathrm{i}\tau^3)^{i_y}$。）求 $\chi$ 与 $\chi_1$ 之间的关系。
\item 证明当 $\chi$ 与 $\chi_1$ 以该方式相联系时，$\pi$ 通量态与 \emph{d} 波态具有相同的自旋子谱。
\end{enumerate}

\subsection{$SU(2)$ 规范涨落的动力学}

\begin{keypoints}
\item Anderson--Higgs 机制可以通过凝聚 $SU(2)$ 规范通量来实现，而无需使用希格斯玻色子。
\item 共线的 $SU(2)$ 通量把 $SU(2)$ 规范结构破缺为 $U(1)$ 规范结构。
% 【块界备注】书页 379（p394）止于此。上述 keypoints 列表的第三条要点印于书页 380（p395）顶部：
%   "Non-collinear SU(2) flux breaks the SU(2) gauge structure down to a Z2 gauge structure.
%    In this case, all of the SU(2) gauge bosons gain a gap."
%   建议译文："非共线的 $SU(2)$ 通量把 $SU(2)$ 规范结构破缺为 $Z_2$ 规范结构。此时，所有的
%   $SU(2)$ 规范玻色子都获得能隙。"
%   该条要点与其后的正文（"Just like the $U(1)$ projective construction, ..."）均属 ch9_c
%   （书页 380 起）。请 ch9_c 在其文件头 %--B-TAIL-- 区或正文开头另起 keypoints 环境续排
%   该条要点（本块环境已闭合，勿并入），勿弃勿重。

% ------------------ 新增术语（ch9_b，供术语表.md "新增术语"区追加） ------------------
% | 英文 | 中文 |
% |---|---|
% | chiral spin state | 手征自旋态 |
% | chiral spin liquid | 手征自旋液体 |
% | semion | 半子 |
% | slave fermion | 从属费米子 |
% | valence band | 价带 |
% | conduction band | 导带（ch5_a 已收） |
% | Landau band | 朗道能带 |
% | (un)frustrated ansatz | （非）阻挫拟设 |
% | Lagrangian multiplier | 拉格朗日乘子 |
% | commensuration | 公度 |
% | pi-flux phase / state | $\pi$ 通量相 / $\pi$ 通量态 |
% | d-wave state | d 波态 |
% | doublet | 二重态 |
% | dressed spinon | 缀饰自旋子 |
% | flux tube | 通量管（ch7_a 的"磁通管"为磁场情形） |
% | Anderson–Higgs mechanism | Anderson–Higgs 机制 |
% | over-counting problem | 重复计数问题 |
% | Dirac algebra | 狄拉克代数 |
% | massless Dirac fermion | 无质量狄拉克费米子 |
% | traceless hermitian matrix | 无迹厄米矩阵 |
